\section{پاسخ سوال ۲}
\begin{stepbox}
\field{محاسبه ابعاد تصویر یک نقطه نورانی در بی‌نهایت}{
ابتدا شرایط دوربین را مدل‌سازی می‌کنیم:
\begin{itemize}
    \item فاصله کانونی لنز: $f = 50 \text{ mm}$
    \item قطر دریچه (Aperture): $D = 26 \text{ mm}$
    \item فاصله جسمی که دوربین روی آن فوکوس شده: $d_o = 2 \text{ m} = 2000 \text{ mm}$
    \item ابعاد هر پیکسل سنسور: $4 \, \mu\text{m} \times 4 \, \mu\text{m}$
\end{itemize}

\textbf{گام اول: محاسبه فاصله سنسور تا لنز ($d_i$)}
دوربین روی فاصله ۲۰۰۰ میلی‌متری فوکوس است. با استفاده از فرمول لنز نازک:
\[ \frac{1}{d_o} + \frac{1}{d_i} = \frac{1}{f} \implies \frac{1}{2000} + \frac{1}{d_i} = \frac{1}{50} \]
\[ \frac{1}{d_i} = \frac{1}{50} - \frac{1}{2000} = \frac{40 - 1}{2000} = \frac{39}{2000} \implies d_i = \frac{2000}{39} \approx 51.282 \text{ mm} \]
بنابراین پرده تصویر (سنسور) در فاصله ۵۱.۲۸۲ میلی‌متری از لنز قرار دارد.

\textbf{گام دوم: بررسی نقطه نورانی در دور دست}
یک نقطه نورانی در بی‌نهایت ($D_o = \infty$) پرتوهای موازی به سمت لنز می‌فرستد. این پرتوها دقیقاً در فاصله کانونی یعنی $d_{\text{focus}} = f = 50 \text{ mm}$ متمرکز می‌شوند.
اما از آنجا که سنسور در فاصله $d_i = 51.282 \text{ mm}$ قرار دارد، پرتوها پس از فوکوس شدن دوباره از هم باز شده و روی سنسور یک لکه دایره‌ای (Circle of Confusion) ایجاد می‌کنند.

\textbf{گام سوم: محاسبه قطر دایره اغتشاش ($c$)}
با استفاده از تشابه مثلث‌ها بین لنز و نقطه فوکوس، و نقطه فوکوس و سنسور:
\[ \frac{c}{d_i - f} = \frac{D}{f} \implies c = D \times \frac{d_i - f}{f} \]
\[ c = 26 \times \frac{\frac{2000}{39} - 50}{50} = 26 \times \frac{2000 - 1950}{39 \times 50} = 26 \times \frac{50}{1950} = \frac{26}{39} = \frac{2}{3} \text{ mm} \approx 0.6667 \text{ mm} = 666.67 \, \mu\text{m} \]

\textbf{گام چهارم: محاسبه تعداد پیکسل‌ها}
این نقطه نورانی به شکل یک دایره روشن با قطر ۶۶۶.۶۷ میکرومتر نمایان می‌شود. مساحت این دایره برابر است با:
\[ Area = \pi \left(\frac{c}{2}\right)^2 = 3.1415 \times (333.33)^2 \approx 349,065.85 \, \mu\text{m}^2 \]
از طرفی مساحت هر پیکسل $4 \times 4 = 16 \, \mu\text{m}^2$ است. تعداد پیکسل‌های پوشش داده شده برابر است با:
\[ \text{تعداد پیکسل} = \frac{349065.85}{16} \approx 21816.6 \text{ پیکسل} \]
لذا این نقطه به شکل یک دایره مات متشکل از حدود \textbf{۲۱۸۱۷ پیکسل} در تصویر ثبت می‌شود.
}
\end{stepbox}
